\chapter{\ifenglish System Workflow\else การทำงานของระบบ\fi}

\section{เครื่องมือคัดกรองสุขภาพช่องปาก}

\begin{tcolorbox}[bluebox]
  \textbf{Oral Health Assessment Tool (OHAT)} ---
  ประเมินสภาพช่องปากเพื่อคัดกรองความเสี่ยง
  และตัดสินใจว่าควรส่งภาพเข้า AI วิเคราะห์หรือไม่

  \vspace{0.5em}

  \textbf{Oral Health Impact Profile (OHIP)} ---
  ประเมินผลกระทบของสุขภาพช่องปากต่อคุณภาพชีวิต
  ทั้งความเจ็บปวด การใช้ชีวิตประจำวัน และการเข้าสังคม
\end{tcolorbox}

\section{กระบวนการยืนยันตัวตนและ Flow การทำงาน}

\subsection{ขั้นตอนการยืนยันตัวตน}

\begin{enumerate}[leftmargin=1.8em, itemsep=3pt]
  \item ผู้ใช้งานกรอกหมายเลขโทรศัพท์
  \item ระบบตรวจสอบสถานะการลงทะเบียนในทะเบียนราษฎร์
  \item \textit{กรณีผู้ใช้ใหม่:} ลงทะเบียน Consent ตาม PDPA และยืนยันบทบาท
  \item ส่งและยืนยัน OTP ผ่าน SMS
  \item ให้สิทธิ์ตามบทบาท (ผู้ป่วย หรือ อสม./ผู้คัดกรอง)
\end{enumerate}

%% >> แทนที่ด้วย: fig_auth_flow.png
%% ที่มา: สไลด์ที่ 12
\begin{figure}[H]
  \centering
  \placeholder{fig\_auth\_flow.png}{Diagram Flow การยืนยันตัวตน พร้อมภาพหน้าจอ Mobile (สไลด์ 12)}
  \caption{ขั้นตอนการยืนยันตัวตนของ eDENT-ELDER:
    กรอกเบอร์โทร ยืนยัน OTP ลง Consent PDPA และได้รับสิทธิ์ตามบทบาท}
\end{figure}

\subsection{Flow การทำงานระบบแบบ 5 เลน}

\renewcommand{\arraystretch}{1.4}
\begin{tabularx}{\linewidth}{|l|X|}
  \hline
  \rowcolor{primaryblue}
  \textcolor{white}{\textbf{เลน}} & \textcolor{white}{\textbf{หน้าที่}}\\
  \hline
  \textbf{เลน 1} \textemdash{} ผู้ป่วย / อสม. &
    รับข้อมูลและ Consent; แบบสอบถาม OHAT/OHIP;
    ถ่ายภาพผ่าน Quality Gate; ส่งเข้า AI\\
  \hline
  \textbf{เลน 2} \textemdash{} ระบบ AI &
    ตรวจสอบคุณภาพภาพ; ตรวจจับรอยโรคและประเมินความเสี่ยง;
    จัดการ Fallback\\
  \hline
  \textbf{เลน 3} \textemdash{} ผู้คัดกรอง &
    Screener Dashboard; Triage ผล AI;
    Flag Retrain หรือ Escalate\\
  \hline
  \textbf{เลน 4} \textemdash{} ผู้เชี่ยวชาญ &
    ยืนยันวินิจฉัย; แจ้งเตือน LINE; ติดตาม 14 วัน; ปิดเคส\\
  \hline
  \textbf{เลน 5} \textemdash{} ผู้ดูแลระบบ / วิชาการ &
    Dashboard ระดับประชากร; แผนที่ความเสี่ยงจังหวัด; กำกับระบบ\\
  \hline
\end{tabularx}

\vspace{8pt}

%% >> แทนที่ด้วย: fig_system_flow.png
%% ที่มา: สไลด์ที่ 13
\begin{figure}[H]
  \centering
  \placeholder{fig\_system\_flow.png}{Diagram Flow 5 เลนพร้อมภาพหน้าจอทุกกลุ่มผู้ใช้ (สไลด์ 13)}
  \caption{Flow การทำงานระบบ eDENT-ELDER แบบ 5 เลน}
\end{figure}

\subsection{อินเทอร์เฟซ Mobile สำหรับผู้ป่วยและ อสม.}

%% >> แทนที่ด้วย: fig_patient_vhv_ui.png
%% ที่มา: สไลด์ที่ 14
\begin{figure}[H]
  \centering
  \placeholder{fig\_patient\_vhv\_ui.png}{หน้าจอ Mobile ผู้ป่วยและ อสม. (สไลด์ 14)}
  \caption{อินเทอร์เฟซ Mobile: รายชื่อผู้ป่วย ผล OHAT แบบ OHIP
    การถ่ายภาพ AI-Guided และหน้าแสดงผล AI (จากซ้ายไปขวา)}
\end{figure}

\subsection{Screener Dashboard}

%% >> แทนที่ด้วย: fig_screener_dashboard.png
%% ที่มา: สไลด์ที่ 15
\begin{figure}[H]
  \centering
  \placeholder{fig\_screener\_dashboard.png}{Screener Dashboard บน Tablet (สไลด์ 15)}
  \caption{Screener Dashboard: คิวภาพเคสพร้อม AI Tag (ซ้าย)
    รายละเอียดเคสพร้อม Annotation (กลาง) และตาราง OHAT/OHIP (ขวา)}
\end{figure}

\subsection{Specialist Dashboard}

%% >> แทนที่ด้วย: fig_specialist_dashboard.png
%% ที่มา: สไลด์ที่ 16
\begin{figure}[H]
  \centering
  \placeholder{fig\_specialist\_dashboard.png}{Specialist Dashboard บน Tablet (สไลด์ 16)}
  \caption{Specialist Dashboard: คิว Referral (ซ้าย) รายละเอียดเคส (กลาง)
    แผนที่ความเสี่ยงระดับจังหวัด (ขวา)}
\end{figure}

\subsection{ภาพหน้าจอจากโปสเตอร์โครงงาน}

%% >> แทนที่ด้วย: fig_poster_interface.png
%% ที่มา: โปสเตอร์ v1 คอลัมน์ขวา บน
\begin{figure}[H]
  \centering
  \placeholder{fig\_poster\_interface.png}{อินเทอร์เฟซแพลตฟอร์มจากโปสเตอร์ v1 (บน)}
  \caption{eDENT-ELDER Platform Interface (โปสเตอร์ v1):
    ระบบ AI ช่วยตรวจจับรอยโรค แสดง Workflow การคัดกรองแบบ Mobile-first}
\end{figure}

%% >> แทนที่ด้วย: fig_poster_screener.png
%% ที่มา: โปสเตอร์ v1 คอลัมน์ขวา ล่าง
\begin{figure}[H]
  \centering
  \placeholder{fig\_poster\_screener.png}{หน้าจอผู้เชี่ยวชาญและผู้คัดกรองจากโปสเตอร์ v1 (ล่าง)}
  \caption{อินเทอร์เฟซวินิจฉัยสำหรับผู้เชี่ยวชาญ (บน) และผู้คัดกรอง (ล่าง)
    พร้อมการเชื่อมต่อผ่าน QR Code}
\end{figure}

